\begin{figure}[h]
    \centering
    \includegraphics[width=0.75\linewidth]{SmartSelect_20250221_112434_Samsung Notes.jpg}
    \label{fig:enter-label}
\end{figure}\\

Para una lente delgada en aire (\( n_0 = 1 \)), la ecuación del fabricante de lentes es:

\[
\frac{1}{f} = (n - 1) \left( \frac{1}{R_1} - \frac{1}{R_2} \right).
\]

Dado que la lente es biconvexa y ambos radios de curvatura son iguales, tomamos \( R_1 = R \) (donde R será negativo pues da la cara "chipoteada" hacia nosotros respecto a un rayo de luz entrante por la izquierda) y \( R_2 = -R \) (lo contrario). Sustituyéndolos en la ecuación:

\[
\frac{1}{f} = (n - 1) \left( \frac{1}{-R} - \frac{1}{R} \right).
\]

\[
\frac{1}{f} = (n - 1) \left( -\frac{2}{R} \right).
\]

\[
R = -2(n - 1) f.
\]

Sustituyendo \( n = 1.5 \) y \( f = 15.0 \) cm:

\[
R = -2(1.5 - 1)(15.0) = -2(0.5)(15.0) = -15.0 \text{ cm}.
\]

Por lo tanto, los radios de curvatura de la lente son:

\[
R_1 = -15.0 \text{ cm}, \quad R_2 = 15.0 \text{ cm}.
\]

Utilizamos la ecuación de lentes delgadas:

\[
\frac{1}{f} = \frac{1}{s} + \frac{1}{s_i}.
\]

Nos dan la posición del objeto \( s = 2.0 \) cm, y queremos demostrar que la imagen se forma en \( s_i = -2.3 \) cm. Sustituyamos los valores:

\[
\frac{1}{15.0} = \frac{1}{2.0} + \frac{1}{s_i}.
\]

Calculamos \( \frac{1}{2.0} \):

\[
\frac{1}{15.0} = 0.0667, \quad \frac{1}{2.0} = 0.5.
\]

Despejamos \( \frac{1}{s_i} \)

\[
\frac{1}{s_i} = \frac{1}{15.0} - \frac{1}{2.0} = 0.0667 - 0.5 = -0.4333.
\]

\[
s_i = \frac{1}{-0.4333} \approx -2.3 \text{ cm}.
\]

Por lo tanto, la imagen se forma a \( s_i = -2.3 \) cm, lo que confirma el resultado esperado.

1. ¿Es real o virtual? 
   Como \( s_i \) es negativo, la imagen se encuentra en el lado del objeto. Esto indica que la imagen es virtual.

2. ¿Está invertida o derecha?
   Derecha, pues la imagen no está invertida y $h_i$ es positiva

3. ¿Qué tan grande es respecto al objeto?  
   La magnificación de la imagen se encuentra con:

   \[
   M = \frac{-s_i}{s} = \frac{-(-2.3)}{2.0} = \frac{2.3}{2.0} = 1.15.
   \]

   Como \( M > 1 \), la imagen es más grande que el objeto por un factor de 1.15 veces.

\begin{figure}[h]
    \centering
    \includegraphics[width=1\linewidth]{traz.png}
\end{figure}